% 16-boundaries.tex: чем Полярис намеренно отказывается становиться.
\section{Чем Полярис отказывается становиться}
\label{sec:boundaries}

\motif{Приватность не есть свобода от принуждения.}

Мощная идентификационная инфраструктура несёт риски, которых не устраняет никакая правильность,
потому что это риски того, для чего правильную систему можно употребить. Репозиторий называет их
ограничениями проектирования и управления, а не политикой, и настоящий раздел излагает их так же.
Каждый есть способ, которым работающий Полярис можно было бы превратить в то, что запрещает его
конституция, и каждый изложен не сильнее, чем позволяют имеющиеся свидетельства.

\begin{description}
  \item[Личность не должна становиться деньгами.] Удостоверение, несущее платёжные полномочия,
    наследует программируемость, и всякое ограничение, какое кому-либо хочется наложить на траты,
    нарастает на идентификационный токен, покуда системе нельзя будет указать, кому отказать у
    бензоколонки. С10 запрещает денежное поле в схеме; если когда-нибудь будет построена система
    ценностей, она живёт в отдельной базе данных, которая потребляет доказательства проверки через
    границу, и несущей является именно граница. Это самое весомое архитектурное решение во всём
    репозитории, и единственное, срок которого не истекает ни с каким этапом.
  \item[Приватность не есть одна лишь тайна.] Система может знать о человеке очень мало и всё же
    принуждать, если каждая дверь требует её разрешения. Нулевое разглашение по умолчанию
    ограничивает то, что записывается; оно не ограничивает того, чего требует общество. Ради этого и
    существует слой многоуровневой регистрации: освобождение есть полноправное состояние в одну
    строку, незарегистрированных нельзя перечислить по случайности устройства схемы, а человека без
    документов можно зарегистрировать через поручителя, чьё поручительство не оставляет на
    удостоверении никакой метки. План сосуществования называет закат старого удостоверения тем
    моментом, когда новое становится обязательным, и отказывается вычислять этот вердикт из того,
    что видит эмитент. Собственная модель угроз репозитория называет сильнейшей атакой на этот слой
    превращение отказа от регистрации в гражданскую невозможность снаружи (<<нет токена, нет
    медицинской помощи>>) и прямо говорит, что схема этого остановить не может, а может лишь
    сделать это поступком, у которого есть имя.
  \item[Противодействие принуждению есть призвание, а не утверждение.] Призвание над десятью
    ограничениями гласит, что никого нельзя заставить сдать свою личность. Механизм, который ему
    служит, код принуждения, противостоит тому принуждающему, ради которого построен, и не
    противостоит двум другим, названным операционным договором, а против законного или
    ведомственного доступа его неизгладимая запись становится свидетельством против самого
    держателя (раздел~\ref{sec:core}). Поэтому парадная дверь говорит \emph{осведомлённый о
    принуждении}, а проверка отвергает \emph{устойчивый к принуждению}. Граница проведена честно:
    система замечает принуждение, но не противостоит ему, и призвание, направленное против того
    противника, перед которым механизм слабее всего, было бы противоположностью того, что обещало
    слово.
  \item[Техническая федерация не есть общественная децентрализация.] Полярис федеративен в коде:
    у каждого органа свой ключ, нет центрального хранилища, нет транзитивного доверия, а
    оператор привязан к одному органу даже внутри общего экземпляра. Два экземпляра,
    запущенные одним автором на одной машине, показывают протокол и ничего не говорят о том, кто
    его эксплуатирует. Независимы ли на деле органы, свидетели и доверяющие стороны, есть свойство
    развёртывания, а не репозитория, и настоящий документ помечает всякое такое утверждение как
    требующее внешней проверки.
  \item[Отзыв не должен становиться гражданским стиранием.] Отозванный токен есть конечное
    состояние в записи, работающей только на добавление, а не удаление; утраченные токены остаются
    утраченными, и их данные не стираются. Преемство идёт по ссылке. Массовый отзыв ограничен
    потолком и соподписантом, а сам потолок можно изменить лишь с записанной причиной. Стирание,
    когда его требует право какой-либо страны, есть замена имени псевдонимом с записью только на
    добавление о том, кто, когда и почему это сделал, но никогда о прежнем имени; сворачивание
    пилота есть именно такое стирание, исполненное, а не обещанное. Орган может прекратить
    удостоверение человека; схема отказывается позволить ему прекратить историю этого человека.
  \item[Соотносимость у доверяющих сторон не то же самое, что несвязываемость у эмитента.] База
    данных эмитента не способна восстановить граф проверок с нулевым разглашением. Две доверяющие
    стороны способны соединить свои журналы по любому устойчивому материалу, который им показало
    предъявление. Оба предложения истинны одновременно, и версия 3 может сказать, какая доля второго
    измерена: на родном пути с нулевым разглашением сговорившаяся пара не имеет преимущества сверх
    множества анонимности против противника, сопоставляющего по равенству, и ничто не задаёт этому
    множеству нижнего предела выше единицы; на пути совместимости одно предъявление оставляет девять
    подписанных эмитентом значений одинаковыми у всякого верификатора (раздел~\ref{sec:privacy}).
    Читатель, который слышит слово \emph{несвязываемый} и представляет любой из путей более сильным,
    введён в заблуждение, и поэтому репозиторий излагает границу утвердительно, а верификатор
    сообщает, какого рода предъявление ему показали.
  \item[Совместимость приносит с собой свойства экосистемы.] Разговор на том формате, на котором
    говорит сторонний кошелёк, был первым, чем воспользовался кто-то извне, и этот формат принёс с
    собой свою связываемость. Репозиторий записывает это как цену совместимости, а не как повод её
    прекратить, и не позволяет пути совместимости заимствовать утверждения о приватности у родного
    пути.
  \item[Никакого графа населения.] Атлас ограничен и никогда не определяет местоположения события
    с нулевым разглашением; Афине запрещены узел-человек, устойчивый суррогат субъекта и обход по
    субъектам; схема не несёт ни одного признака, по которому можно было бы отфильтровать население;
    квоты и оповещения привязаны к ведомству и никогда к человеку; отчёт о прозрачности не несёт
    итога, из которого можно было бы восстановить малое значение. Отказ структурен, потому что граф
    населения, однажды построенный, переживает намерения того, кто его строил.
  \item[Машинное обучение не должно превращать идентификационную инфраструктуру в систему оценок.]
    Предложенный обучающийся слой вправе изучать, как ведёт себя идентификационная система, и
    никогда не вправе изучать, как ведут себя люди. Оценки мошенничества или риска для отдельных
    людей, поведенческие профили, выведение связей между людьми, оценки права на что-либо или
    вероятности отзыва, кластеризация и векторные представления на уровне человека и отказ в правах
    по решению модели запрещены по умолчанию, а рекомендация есть свидетельство, но никогда не
    полномочие.
\end{description}

Репозиторий перечисляет также постоянные <<не-цели>>, срок которых не истекает ни с каким этапом
дорожной карты: никаких платёжных рельсов, никакой центральной биометрической базы, никакой
фильтрации и аналитики признаков в масштабе населения, никаких социальных рейтингов, никакого
коммерческого или рекламного использования данных о проверках и никаких настоящих личных данных,
пока не существует рамки согласия. Каждое из этого есть способ сделать систему полезнее для
кого-то, и каждое отвергается с первого взгляда. Со времён операционного договора ещё один отказ
управляет самим репозиторием: никакая новая возможность не строится потому, что была бы изящной,
полной или желанной; изменению продукта нужна названная внешняя сторона, которой оно требуется,
либо исполнимый контрпример против того, что уже обещано.

\analogy{Законы над городом связывают сам город. Ратуша не вправе чеканить монету. Реестр нельзя
превратить в список тех, кого нет. Ворота нельзя закрыть перед соседом по слову соседа этого
соседа. Имя, вычеркнутое из реестра, остаётся читаемым как вычеркнутое. Ни одного писаря, каким бы
толковым он ни был, нельзя посадить угадывать, кто из горожан, скорее всего, наделает бед. И город
не утверждает, что его стены защитят человека от стражника, который стоит с ним рядом.}
